% ===== PRÉAMBULE CANONIQUE — à coller VERBATIM en tête de CHAQUE .tex =====
\documentclass[11pt,a4paper]{article}
\usepackage[utf8]{inputenc}\usepackage[T1]{fontenc}\usepackage{lmodern}
\usepackage[french]{babel}
\usepackage{amsmath,amssymb,mathtools}\usepackage{icomma}
\usepackage{siunitx}\sisetup{locale=FR}  % \qty{}{}, \si{}, \ang{}
\usepackage{xcolor}\usepackage{colortbl}
\usepackage{tikz}\usepackage{pgfplots}\pgfplotsset{compat=1.18}
\usepackage[siunitx,europeanresistors]{circuitikz} % circuits (élec.)
\usepackage[most]{tcolorbox}\usepackage{enumitem}\usepackage{array,booktabs}
\usepackage[a4paper,margin=2.2cm]{geometry}
\usetikzlibrary{arrows.meta,calc,angles,quotes,positioning,patterns,%
  backgrounds,shapes.geometric,decorations.pathmorphing,decorations.markings,intersections}
\definecolor{primary}{RGB}{222,49,99}\definecolor{primarydark}{RGB}{150,22,55}
\definecolor{defcol}{RGB}{0,114,178}\definecolor{methcol}{RGB}{52,92,148}
\definecolor{excol}{RGB}{0,135,90}\definecolor{attncol}{RGB}{200,40,55}
\tcbset{boxbase/.style={enhanced,breakable,boxrule=0.9pt,arc=2.5pt,
  left=3mm,right=3mm,top=2.3mm,bottom=2.3mm,fonttitle=\bfseries,coltitle=white}}
% --- boîtes TUTORIEL (noms EXACTS, ne pas renommer) ---
\newtcolorbox{cle}[1][]{boxbase,colback=defcol!6,colframe=defcol,
  title={\textsc{À retenir}\ifx\relax#1\relax\relax\else\textmd{ : \itshape #1}\fi}}
\newtcolorbox{methode}[1][]{boxbase,colback=methcol!7,colframe=methcol,sharp corners,
  title={\textsc{Méthode}\ifx\relax#1\relax\relax\else\textmd{ --- \itshape #1}\fi}}
\newtcolorbox{exemplebox}[1][]{boxbase,colback=excol!5,colframe=excol,
  title={\textsc{Exemple}\ifx\relax#1\relax\relax\else\textmd{ : \itshape #1}\fi}}
\newtcolorbox{piege}[1][]{boxbase,colback=attncol!6,colframe=attncol,
  title={\textsc{Piège}\ifx\relax#1\relax\relax\else\textmd{ : \itshape #1}\fi}}
\newtcolorbox{remarque}[1][]{boxbase,colback=black!4,colframe=black!45,
  title={\textsc{Remarque}\ifx\relax#1\relax\relax\else\textmd{ : \itshape #1}\fi}}
% --- boîtes EXERCICE / CORRIGÉ (noms EXACTS, auto-comptées) ---
\newtcolorbox[auto counter]{exo}[1][]{boxbase,colback=primary!4,colframe=primary,
  title={\textsc{Exercice}~\thetcbcounter\ifx\relax#1\relax\relax\else\textmd{ --- \itshape #1}\fi}}
\newtcolorbox[auto counter]{corrige}[1][]{boxbase,colback=excol!4,colframe=excol,
  title={\textsc{Corrigé}~\thetcbcounter\ifx\relax#1\relax\relax\else\textmd{ --- \itshape #1}\fi}}
\newcommand{\vv}[1]{\vec{#1}}\newcommand{\ds}{\displaystyle}
\newcommand{\R}{\mathbb{R}}\newcommand{\N}{\mathbb{N}}\newcommand{\Z}{\mathbb{Z}}
% (un \usepackage{fancyhdr}+en-tête est possible ; il est ignoré côté web)
% =========================================================================
\begin{document}

\begin{center}
{\large\bfseries\color{excol} Thème 2 — Corrigé Facile}\\[-2pt]
{\color{excol}\rule{5cm}{1pt}}
\end{center}

\begin{corrige}[calculer la pulsation $\omega$]
\begin{enumerate}[label=\alph*)]
  \item $\omega=2\pi f=2\pi\times 10=20\pi\approx\SI{62.8}{\radian\per\second}.$
  \item $\omega=\dfrac{2\pi}{T}=\dfrac{2\pi}{0{,}5}=4\pi\approx\SI{12.6}{\radian\per\second}.$
\end{enumerate}
\end{corrige}

\begin{corrige}[calculer le nombre d'onde $k$]
\begin{enumerate}[label=\alph*)]
  \item $k=\dfrac{2\pi}{\lambda}=\dfrac{2\pi}{0{,}5}=4\pi\approx\SI{12.6}{\radian\per\metre}.$
  \item $k=\dfrac{2\pi}{0{,}2}=10\pi\approx\SI{31.4}{\radian\per\metre}.$
\end{enumerate}
\end{corrige}

\begin{corrige}[lire les paramètres dans l'équation]
On identifie terme à terme avec $y=A\sin(\omega t - kx)$.
\begin{enumerate}[label=\alph*)]
  \item $A=\SI{0.05}{\metre}$ (le facteur devant le sinus).
  \item Facteur devant $t$ : $\omega=10\pi\approx\SI{31.4}{\radian\per\second}$ ;
        facteur devant $x$ : $k=4\pi\approx\SI{12.6}{\radian\per\metre}$.
  \item Le terme en $x$ est $-kx$ (signe moins) : l'onde se propage vers les $x$
        \textbf{croissants}.
\end{enumerate}
\end{corrige}

\begin{corrige}[vitesse et accélération maximales d'un point]
\begin{enumerate}[label=\alph*)]
  \item $v_{\max}=A\omega=0{,}02\times 50=\SI{1}{\metre\per\second}.$
  \item $a_{\max}=A\omega^{2}=0{,}02\times 50^{2}=0{,}02\times 2500=\SI{50}{\metre\per\second\squared}.$
\end{enumerate}
\end{corrige}

\begin{corrige}[concordance ou opposition de phase ?]
On mesure la distance entre les points puis on la compare à $\lambda=\SI{4}{\metre}$ :
$d=n\lambda\Rightarrow$ concordance ; $d=(2n+1)\lambda/2\Rightarrow$ opposition.
\begin{enumerate}[label=\alph*)]
  \item $M$ ($x=1$) et $P$ ($x=5$) : $d=\SI{4}{\metre}=\lambda$ $\Rightarrow$
        \textbf{concordance} (les deux sont sur une crête).
  \item $M$ ($x=1$) et $N$ ($x=3$) : $d=\SI{2}{\metre}=\lambda/2$ $\Rightarrow$
        \textbf{opposition} ($M$ sur une crête, $N$ dans un creux).
\end{enumerate}
\end{corrige}

\end{document}
